%%%%%%%%%%%%%%%%%%%%%%%%%%%%%%%%%%%%%%%%%%%%%%%%%%%%%%%%%%%%%%%%%%%%%%%%%%%
\chapter{Tokenization}
\label{ch:tokenization}
%%%%%%%%%%%%%%%%%%%%%%%%%%%%%%%%%%%%%%%%%%%%%%%%%%%%%%%%%%%%%%%%%%%%%%%%%%%

Tokenization is the process of splitting text into the discrete units a model processes.
Understanding it is essential for working with embedding models because token counts---not
character counts---determine whether a text fits in a model's context window, how chunking
interacts with model limits, and why costs and latency vary with input length.

%%%%%%%%%%%%%%%%%%%%%%%%%%%%%%%%%%%%%%%%%%%%%%%%%%%%%%%%%%%%%%%%%%%%%%%%%%%
\section{Why Tokens, Not Words or Characters?}
%%%%%%%%%%%%%%%%%%%%%%%%%%%%%%%%%%%%%%%%%%%%%%%%%%%%%%%%%%%%%%%%%%%%%%%%%%%

Early NLP systems split on whitespace: each word is a unit. This breaks down on several
practical problems:

\begin{itemize}
    \item \textbf{Vocabulary explosion.} English has hundreds of thousands of words, and new
          ones are coined constantly. A word-level vocabulary must either be huge or mark
          unknown words as \texttt{[UNK]}, losing all information about them.
    \item \textbf{Morphological variants.} ``run'', ``runs'', ``running'', and ``ran'' are four
          different vocabulary entries that share almost all meaning. A word-level model treats
          them as unrelated.
    \item \textbf{Multilinguality.} Japanese and Chinese have no spaces. Arabic and Hebrew write
          without vowels. A word tokenizer that splits on spaces fails silently.
\end{itemize}

Character-level tokenization solves the vocabulary problem (26 letters plus punctuation is a
tiny vocabulary) but produces very long sequences. Attention is O($n^2$) in sequence length,
so character sequences are expensive to process.

\textbf{Subword tokenization}---the dominant approach---sits between the two. Common words stay
as single tokens. Rare words are split into meaningful subword pieces. Sequences stay short
enough for efficient attention.

%%%%%%%%%%%%%%%%%%%%%%%%%%%%%%%%%%%%%%%%%%%%%%%%%%%%%%%%%%%%%%%%%%%%%%%%%%%
\section{BPE (Byte Pair Encoding)}
%%%%%%%%%%%%%%%%%%%%%%%%%%%%%%%%%%%%%%%%%%%%%%%%%%%%%%%%%%%%%%%%%%%%%%%%%%%

BPE is the most widely used tokenization algorithm (GPT-2 through GPT-4o, LLaMA, Qwen, Mistral).

\textbf{How it works:}
\begin{enumerate}
    \item Start with a vocabulary of individual characters (or bytes).
    \item Count all adjacent pairs in the training corpus.
    \item Merge the most frequent pair into a new token.
    \item Repeat until the vocabulary reaches the target size (typically 32k--128k tokens).
\end{enumerate}

\textbf{Example:}

\begin{Verbatim}
Training text: "low lower lowest lower low"

Initial: l o w  l o w e r  l o w e s t  l o w e r  l o w
Step 1: merge 'l'+'o' → 'lo' (most frequent pair)
        lo w  lo w e r  lo w e s t  lo w e r  lo w
Step 2: merge 'lo'+'w' → 'low'
        low  low e r  low e s t  low e r  low
Step 3: merge 'low'+'e' → 'lowe'
        low  lowe r  lowe s t  lowe r  low
\end{Verbatim}

After training on a large corpus, common words like ``the'', ``is'', and ``model'' are single
tokens. Rare words split into recognizable subwords: ``tokenization'' becomes
\texttt{["token", "ization"]}. Truly novel words split further:
\texttt{"Qwen3Embedding"} $\to$ \texttt{["Qw", "en", "3", "Embed", "ding"]}.

The \textbf{byte-level variant} (used by GPT-2 and descendants) operates on raw bytes rather
than characters. This ensures any Unicode text can be tokenized without an unknown-token
fallback, at the cost of longer sequences for non-Latin scripts.

%%%%%%%%%%%%%%%%%%%%%%%%%%%%%%%%%%%%%%%%%%%%%%%%%%%%%%%%%%%%%%%%%%%%%%%%%%%
\section{WordPiece}
%%%%%%%%%%%%%%%%%%%%%%%%%%%%%%%%%%%%%%%%%%%%%%%%%%%%%%%%%%%%%%%%%%%%%%%%%%%

Used by BERT and its descendants (including many embedding models like BGE and
nomic-embed-text). Conceptually similar to BPE, but merges are chosen to maximize the
likelihood of the training data under a language model, not just by raw pair frequency.
Subword pieces are prefixed with \texttt{\#\#} to indicate continuation:

\begin{Verbatim}
"playing"     → ["play", "##ing"]
"unplayable"  → ["un", "##play", "##able"]
\end{Verbatim}

%%%%%%%%%%%%%%%%%%%%%%%%%%%%%%%%%%%%%%%%%%%%%%%%%%%%%%%%%%%%%%%%%%%%%%%%%%%
\section{SentencePiece}
%%%%%%%%%%%%%%%%%%%%%%%%%%%%%%%%%%%%%%%%%%%%%%%%%%%%%%%%%%%%%%%%%%%%%%%%%%%

Used by LLaMA, Mistral, and many multilingual models. Operates directly on raw text without
pre-tokenization---no whitespace splitting required first. Handles languages without spaces
naturally. The underlying algorithm is typically BPE or Unigram.

%%%%%%%%%%%%%%%%%%%%%%%%%%%%%%%%%%%%%%%%%%%%%%%%%%%%%%%%%%%%%%%%%%%%%%%%%%%
\section{Token Count vs Character Count}
%%%%%%%%%%%%%%%%%%%%%%%%%%%%%%%%%%%%%%%%%%%%%%%%%%%%%%%%%%%%%%%%%%%%%%%%%%%

A critical practical point: \textbf{model limits are in tokens, not characters or words}.

Approximate conversions for English text:

\begin{center}
\begin{tabular}{ll}
\toprule
\textbf{Unit} & \textbf{Typical ratio} \\
\midrule
1 token & \textasciitilde4 characters \\
1 token & \textasciitilde0.75 words \\
100 tokens & \textasciitilde75 words, \textasciitilde400 characters \\
1{,}000 tokens & \textasciitilde750 words, \textasciitilde3--4 paragraphs \\
\bottomrule
\end{tabular}
\end{center}

These ratios are for English. Other languages differ:

\begin{center}
\begin{tabular}{ll}
\toprule
\textbf{Language} & \textbf{Tokens per word (approx)} \\
\midrule
English & 1.3 \\
Spanish / French & 1.4 \\
German & 1.5 (compound words split more) \\
Chinese & 1.5--2.5 (ideographs often multi-byte) \\
Arabic & 1.5--2.0 \\
Code (Python) & 2--4 (identifiers, indentation, symbols) \\
\bottomrule
\end{tabular}
\end{center}

Code tokenizes much less efficiently than prose. A 100-line Python function may consume
300--500 tokens. This matters for code retrieval use cases where chunk size must be set by
token count, not character count.

%%%%%%%%%%%%%%%%%%%%%%%%%%%%%%%%%%%%%%%%%%%%%%%%%%%%%%%%%%%%%%%%%%%%%%%%%%%
\section{How Token Limits Affect Chunking}
%%%%%%%%%%%%%%%%%%%%%%%%%%%%%%%%%%%%%%%%%%%%%%%%%%%%%%%%%%%%%%%%%%%%%%%%%%%

Every embedding model has a maximum context length:

\begin{center}
\begin{tabular}{ll}
\toprule
\textbf{Model} & \textbf{Max tokens} \\
\midrule
BGE-small-en-v1.5 & 512 \\
nomic-embed-text-v1.5 & 8{,}192 \\
text-embedding-3-small & 8{,}191 \\
text-embedding-3-large & 8{,}191 \\
Qwen3-Embedding series & 32{,}768 \\
\bottomrule
\end{tabular}
\end{center}

What happens when input exceeds the limit depends on the implementation:

\begin{itemize}
    \item \textbf{Truncation:} the model silently discards tokens beyond the limit. The embedding
          represents only the first $N$ tokens. Content in the truncated portion is invisible
          to retrieval.
    \item \textbf{Error:} the API or model raises an error. Callers must pre-split.
\end{itemize}

Truncation is the silent failure mode. If your chunks regularly exceed the model's token limit,
you are losing content without any warning.

Chunk size should be set in tokens when precision matters. For a 512-token limit model with a
15\% safety margin, target 435-token chunks. Converted to characters at 4 chars/token: roughly
1{,}740 characters. memvid's default of 1{,}200 characters is conservative for 512-token models
and well within limits for 8K+ token models.

%%%%%%%%%%%%%%%%%%%%%%%%%%%%%%%%%%%%%%%%%%%%%%%%%%%%%%%%%%%%%%%%%%%%%%%%%%%
\section{Tokens and Embedding Quality}
%%%%%%%%%%%%%%%%%%%%%%%%%%%%%%%%%%%%%%%%%%%%%%%%%%%%%%%%%%%%%%%%%%%%%%%%%%%

Shorter inputs produce different embeddings than longer ones because pooling collapses all
token vectors into one. A 5-token chunk and a 500-token chunk are both represented by a single
1536-dim vector, but the 500-token vector averages over far more content.

Implications:

\begin{itemize}
    \item \textbf{Very short chunks} (1--3 sentences) embed well---the vector captures the
          specific topic of that chunk cleanly.
    \item \textbf{Very long chunks} produce embeddings that represent a blend of all topics in
          the chunk. They retrieve for any of those topics at reduced precision.
    \item \textbf{Asymmetric retrieval:} a 10-word query embedding will not be close to a
          2{,}000-token document embedding even if the document contains the answer, because the
          document's vector is averaged over content unrelated to the query. This is why chunking
          exists.
\end{itemize}

%%%%%%%%%%%%%%%%%%%%%%%%%%%%%%%%%%%%%%%%%%%%%%%%%%%%%%%%%%%%%%%%%%%%%%%%%%%
\section{Special Tokens}
%%%%%%%%%%%%%%%%%%%%%%%%%%%%%%%%%%%%%%%%%%%%%%%%%%%%%%%%%%%%%%%%%%%%%%%%%%%

Most models use reserved tokens for structural purposes:

\begin{center}
\begin{tabular}{lll}
\toprule
\textbf{Token} & \textbf{Meaning} & \textbf{Models} \\
\midrule
\texttt{[CLS]} & Classification / sentence start & BERT-style \\
\texttt{[SEP]} & Separator between segments & BERT-style \\
\texttt{[PAD]} & Padding to uniform length & All \\
\texttt{[UNK]} & Unknown token (rare in modern BPE) & All \\
\texttt{<s>} & Sequence start & LLaMA/Mistral \\
\texttt{</s>} & Sequence end & LLaMA/Mistral \\
\texttt{<|endoftext|>} & End of text & GPT-2/3/4 \\
\bottomrule
\end{tabular}
\end{center}

Special tokens consume context window budget. A BGE model with a 512-token limit uses 2
special tokens (\texttt{[CLS]} and \texttt{[SEP]}), leaving 510 tokens for content.

For CLS pooling models, the \texttt{[CLS]} token is the output used as the embedding.
For last-token pooling models (Qwen3), the final real token before \texttt{</s>} is used.

%%%%%%%%%%%%%%%%%%%%%%%%%%%%%%%%%%%%%%%%%%%%%%%%%%%%%%%%%%%%%%%%%%%%%%%%%%%
\section{Counting Tokens Without Running the Model}
%%%%%%%%%%%%%%%%%%%%%%%%%%%%%%%%%%%%%%%%%%%%%%%%%%%%%%%%%%%%%%%%%%%%%%%%%%%

Counting tokens before sending text to a model is useful for ensuring chunks stay within
limits, estimating API cost, and debugging truncation issues.

\textbf{Python (HuggingFace tokenizers):}

\begin{Verbatim}
from transformers import AutoTokenizer

tokenizer = AutoTokenizer.from_pretrained("BAAI/bge-small-en-v1.5")
tokens = tokenizer.encode("The quick brown fox")
print(len(tokens))  # 6
\end{Verbatim}

\textbf{Python (tiktoken, for OpenAI models):}

\begin{Verbatim}
import tiktoken

enc = tiktoken.get_encoding("cl100k_base")  # text-embedding-3-*
tokens = enc.encode("The quick brown fox")
print(len(tokens))  # 5
\end{Verbatim}

\textbf{Rule of thumb without a tokenizer:}
Divide character count by 4 for English prose. Accurate to $\pm 20\%$ for most inputs.
For code, divide by 3. For non-Latin scripts, divide by 2.

%%%%%%%%%%%%%%%%%%%%%%%%%%%%%%%%%%%%%%%%%%%%%%%%%%%%%%%%%%%%%%%%%%%%%%%%%%%
\section{Practical Takeaways}
%%%%%%%%%%%%%%%%%%%%%%%%%%%%%%%%%%%%%%%%%%%%%%%%%%%%%%%%%%%%%%%%%%%%%%%%%%%

\begin{itemize}
    \item Token limits are per-model and in tokens, not characters. Check the model card.
    \item English prose: \textasciitilde4 characters per token. Code: \textasciitilde3 chars per token. CJK: \textasciitilde2 chars per token.
    \item Exceeding the limit silently truncates in most implementations.
    \item Set chunk size in tokens when working with models near their 512-token limit;
          character estimates are close enough for 8K+ token models.
    \item Short chunks embed their specific topic precisely; long chunks produce averaged embeddings.
    \item Special tokens consume 2--4 tokens of context budget per call.
\end{itemize}
